% Candidate Figure 1 for the staged-v7 solver. Transparent line art only.
\begin{figure}[H]
\centering
\begin{tikzpicture}[
  every node/.style={font=\fontsize{8.4}{9.8}\selectfont},
  box/.style={rectangle,draw=black,line width=0.45pt,rounded corners=0pt,
    align=center,minimum height=8.2mm,inner xsep=2.5pt,inner ysep=1.5pt},
  wide/.style={box,text width=9.1cm},
  stage/.style={box,text width=9.1cm,minimum height=7.6mm},
  view/.style={box,text width=3.15cm,minimum height=9.2mm},
  arrow/.style={-{Latex[length=2.4mm,width=1.55mm]},line width=0.52pt,
    line cap=rect,line join=miter,shorten >=1pt},
  stem/.style={line width=0.52pt,line cap=rect,line join=miter}
]
\node[wide] (input) at (0,0)
  {算例与时空参数\\有限车队、共同可行初解};

\node[stage] (stageone) at (0,-1.05)
  {第一阶段：三视角独立搜索（各5000次）};
\node[view] (s1f) at (-4.15,-2.10)
  {HGS-F\\燃油成本视角};
\node[view] (s1e) at (0,-2.10)
  {HGS-E\\简化电动视角};
\node[view] (s1m) at (4.15,-2.10)
  {HGS-M\\机制感知视角};
\node[wide] (reviewone) at (0,-3.35)
  {分别执行完整模型复核\\建立三份视角精英档案};
\node[wide] (mipone) at (0,-4.35)
  {汇集可行路线；第一次限时MIP组合（30 s）};
\node[wide] (protected) at (0,-5.35)
  {三视角解与组合解择优\\形成阶段1保护解};

\node[stage] (stagetwo) at (0,-6.45)
  {第二阶段：本视角精英延续（各20000次）};
\node[view] (s2f) at (-4.15,-7.50)
  {HGS-F\\同视角热启动};
\node[view] (s2e) at (0,-7.50)
  {HGS-E\\同视角热启动};
\node[view] (s2m) at (4.15,-7.50)
  {HGS-M\\同视角热启动};
\node[wide] (reviewtwo) at (0,-8.75)
  {分别执行完整模型复核\\合并各自的两阶段档案};
\node[wide] (miptwo) at (0,-9.75)
  {扩展可行路线池；第二次限时MIP组合（30 s）};
\node[wide] (judge) at (0,-10.75)
  {完整复算、独立检查与受保护择优};
\node[wide] (output) at (0,-11.75)
  {输出完整模型核验的最好可行解};

\draw[arrow] (input) -- (stageone);
\coordinate (splitone) at (0,-1.58);
\draw[stem] (stageone.south) -- (splitone);
\draw[arrow] (splitone) -| (s1f.north);
\draw[arrow] (splitone) -- (s1e.north);
\draw[arrow] (splitone) -| (s1m.north);
\coordinate (mergeone) at (0,-2.72);
\draw[stem] (s1f.south) |- (mergeone);
\draw[stem] (s1e.south) -- (mergeone);
\draw[stem] (s1m.south) |- (mergeone);
\draw[arrow] (mergeone) -- (reviewone);
\draw[arrow] (reviewone) -- (mipone);
\draw[arrow] (mipone) -- (protected);
\draw[arrow] (protected) -- (stagetwo);

\coordinate (splittwo) at (0,-6.98);
\draw[stem] (stagetwo.south) -- (splittwo);
\draw[arrow] (splittwo) -| (s2f.north);
\draw[arrow] (splittwo) -- (s2e.north);
\draw[arrow] (splittwo) -| (s2m.north);
\coordinate (mergetwo) at (0,-8.12);
\draw[stem] (s2f.south) |- (mergetwo);
\draw[stem] (s2e.south) -- (mergetwo);
\draw[stem] (s2m.south) |- (mergetwo);
\draw[arrow] (mergetwo) -- (reviewtwo);
\draw[arrow] (reviewtwo) -- (miptwo);
\draw[arrow] (miptwo) -- (judge);
\draw[arrow] (judge) -- (output);
\end{tikzpicture}
\caption{两阶段多视角混合遗传搜索流程}
\label{fig:algorithm-flow}
\end{figure}
